\caption{External sequence preferences and source-resolved natural variants. \textbf{a}, All 25 proteins' paired Full/Frequency scores and differences for primary all-word Spearman and secondary composition-matched concordance. Adjacent difference axes share the absolute axes' units and spans, offset to the Frequency means. In the adjacent difference axes, red points and interval bars give frozen paired-protein means and 95\% bootstrap intervals; grey points are individual protein differences. The mean Full-minus-Frequency differences are $+0.146$ (95\% CI $+0.038$ to $+0.248$) for Spearman and $+3.78$ percentage points ($+0.90$ to $+6.73$) for composition-matched concordance. Complete Sequence controls remain in Figure~S31. \textbf{b}, All 30 QKI A$_2$C$_2$U words and 48 adjacent exchanges. Height encodes RNAcompete $Z_5$; node halves encode Full/Sequence mean responses across 32 computational backgrounds on a shared asinh colour scale. Gold marks three CUAAC-neighbour edges. \textbf{c}, Exact ACVR1B A$>$G, source UPF1 preference for A (ASB FDR 0.0446), seven saved model/control responses and published HeLa RNA reporter measurements from \citet{yang2019asb}. Hollow shapes identify three fixed-input permutations, not uncertainty or biological replicates. Reporter points are three biological measurements per allele with mean ticks; RNA abundance and predicted K562 binding are distinct endpoints. This is a post hoc training-locus illustration. \textbf{d}, All 396 exposed ASB-positive events, with source-preferred-minus-other responses and five-arm RBP-equal accuracy/95\% intervals. Red/cyan identify Full-only/Sequence-only direction agreement. \textbf{e}, All eight ClinVar P/LP and/or TCGA-overlap sites, 13 events and four-arm REF-minus-ALT responses. Filled/open hatched bars favour REF/ALT; every source label favours REF. Agreement counts are 7/10/12/9 for Full/Sequence/flat/permuted input; the constant-REF rule gives 13/13. H/K denote HepG2/K562. These events include 11 training and two validation observations, with no strict-test events. Disease annotations do not establish binding truth or clinical causality. Both axes in d and all four effect axes in e share asinh width 0.02 within their panel; ticks retain original log-count units. The reference conservation input is fixed across each allele pair. Exact site identities: Figure~S32. Methods S1.18/S1.20--S1.22; Data S23/S25--S27.}
